The next theorem specifies an integral coefficient order for the generators of \citet[Definition 11.1, equation (93), p.~59]{BSM2013}. It does not assume a positive or canonical basis.

Let $v$ be an indeterminate and put
\[
 D=(v+v^{-1})^2,\qquad S_D=\mathbb Z[D],\qquad
 R_v=\mathbb Z[v,v^{-1}],\qquad
 F_D=\mathbb Q(D),\qquad K_v=\mathbb Q(v).
\]
Thus the earlier Laurent parameter is $q=v^2$. For a positive integer $r$, let $H_{R_v}(\mathfrak S_r)$ be the type-A Hecke algebra with its standard basis $T_w$ and relations
\[
 (T_s-v)(T_s+v^{-1})=0
\]
for the simple reflections, together with the braid relations. In
\[
 M_{R_v}=H_{R_v}(\mathfrak S_r)\otimes_{R_v}
 H_{R_v}(\mathfrak S_r)
\]
define
\[
 P_s=(T_s+v^{-1})\otimes(T_s+v^{-1})
       +(T_s-v)\otimes(T_s-v).
\]
A \emph{word} means a finite product of the $P_s$, including the empty word $1$. Write $A_{S_D},A_{R_v},A_{F_D},A_{K_v}$ for the spans of these words over the indicated rings or fields, all viewed in the ambient Hecke tensor algebra over $K_v$.

\begin{theorem}\label{thm:hecke-word-descent}
For every positive integer $r$, the following statements hold.
\begin{enumerate}
\item The algebra $A_{S_D}$ is a finite torsion-free $S_D$-module.
\item Multiplication induces isomorphisms of algebras
\[
 R_v\otimes_{S_D}A_{S_D}\xrightarrow{\ \sim\ }A_{R_v},
 \qquad
 K_v\otimes_{F_D}A_{F_D}\xrightarrow{\ \sim\ }A_{K_v}.
\]
\item The algebra $A_{F_D}$ has an $F_D$-basis consisting of words. Every multiplication coefficient in such a basis belongs to $F_D$.
\item For $a\in A_{S_D}$, let $L_a$ be left multiplication on the finite-dimensional $F_D$-vector space $A_{F_D}$. Then
\[
 \det(X-L_a)\in S_D[X].
\]
In particular, the regular trace of every word lies in $\mathbb Z[D]$.
\end{enumerate}
\end{theorem}

